\documentclass[11pt,a4paper]{article}
\usepackage[T1]{fontenc}
\usepackage[utf8]{inputenc}
\usepackage[czech]{babel}
\usepackage{amsmath,amssymb,amsthm,mathtools}
\usepackage[margin=2.5cm]{geometry}
\usepackage[hidelinks]{hyperref}
\newtheorem{theorem}{Věta}
\newtheorem{lemma}[theorem]{Lemma}
\newtheorem{proposition}[theorem]{Tvrzení}
\newcommand{\E}{\mathbb E}
\newcommand{\Law}{\operatorname{Law}}
\newcommand{\HH}{\mathsf H}
\title{Posunuté Waringovy ocasy: přesná orbita a zbývající transport hlavy}
\author{}
\date{19. září 2026}
\begin{document}
\maketitle

\noindent\textbf{Rozsah výsledku.}
Níže sestrojujeme explicitní přechod ocasů pro libovolnou historii,
dokazujeme jejich dyadický hmotnostní odhad a přesně charakterizujeme
nezáporné doplnění jednoho fanu. Neprokazujeme, že doplnění lze volit
rekurzivně pro všechny historie. Proto nejde o důkaz
$\mathcal B_n=O(\log n)$ ani o důkaz existence nekonečného exactního stromu.

Položme
\[
 P=(8,16,16,16,8)/64,\qquad Q=(7,20,10,20,7)/64,
 \qquad W_a(j)=\frac{a}{(a+j-1)(a+j)}\quad(j\ge1).
\]
Kořen je $H_2=W_2$. Píšeme $q_d=Q^{*d}$,
$s(h)=\sum_{i=1}^{d}h_i$ pro $|h|=d$ a $\HH_m=\sum_{j=1}^m1/j$.

\section*{1. Přesné rozbalení jedné fáze}
Pro celá $L\ge0$ platí
\begin{equation}\label{peel}
 W_a=\sum_{j=1}^L\frac{a}{(a+j-1)(a+j)}\delta_j
       +\frac{a}{a+L}\,\delta_L*W_{a+L}.
\end{equation}
Je to rovnost koeficientů: za řezem $L$ se čitatel
$\frac{a}{a+L}(a+L)$ rovná $a$. Pro celé $a$ a $L=a$ má pokračující
část hmotnost přesně $1/2$. Tato identita sama ještě neřeší veřejné větvení.

U konečné směsi posunutých ocasů
$\sum_i w_i\delta_{r_i}*W_{a_i}$ definujme kapitálovou míru fází
\[
 \Lambda=\sum_i w_i a_i\,\delta_{a_i-r_i}.
\]
Za všemi jejími řezy je pravděpodobnostní koeficient rezervy $j$ roven
\[
 \int\frac{\Lambda(db)}{(j+b-1)(j+b)}.
\]
Konečná část zákona do tohoto vyjádření nevstupuje.

\begin{lemma}[Orbitální přechod ocasu]\label{phase}
Konvolvujme zdroj s $Q$, zvolme společný řez $M$ nad konečnými částmi
i nad posuny zdroje a rozdělme zdrojovou hmotu nad $M$ mezi veřejné
symboly nezávisle v poměrech $P_x$. Pak kapitálová míra ocasu dítěte $x$ je
\begin{equation}\label{orbit}
 \boxed{\Lambda_x=\sum_{y=0}^4Q_y\,(b\mapsto b+x-y)_\#\Lambda.}
\end{equation}
Zejména $|\Lambda_x|=|\Lambda|$ pro každé dítě.
\end{lemma}
\begin{proof}
Komponenta $w\delta_r*W_a$ po vstupu $y$ a rozbalení v $M$ přispěje
ocasem
\[
 wQ_y\frac{a}{a+M-r-y}\,
 \delta_M*W_{a+M-r-y}.
\]
Po přiřazení symbolu $x$ a podmínění na něj je posun $M-x$ a fáze
$a'=a+M-r-y$. Součin hmotnosti a fáze je $wQ_ya$ a rozdíl fáze a posunu
je $(a-r)+x-y$. To je \eqref{orbit}.
\end{proof}

Pro kořen $W_2$ tak dostáváme pro \emph{každou} konečnou historii
explicitní, nezápornou míru
\begin{equation}\label{lambda}
 \boxed{\Lambda_h=2\,\Law\bigl(2+s(h)-Y_d\bigr),\qquad Y_d\sim q_d.}
\end{equation}
Její podpora může unikat oběma směry. Nejde o dvoufázový ani periodický
model; \eqref{lambda} je přesná formule pro celou orbitální hloubku.
Je to zatím orbita \emph{ocasů}, nikoli již hotových pravděpodobnostních fanů.

\section*{2. Konkrétní řezy a dyadická hmotnost}
Zvolme $K_d=8d$ a $L_h=8d-s(h)$. Hledaný konkrétní tvar zákonů je
\begin{equation}\label{normal}
 B_h=A_h+
 \sum_{y=0}^{4d}\frac{2q_d(y)}{8d+2-y},
       \delta_{8d-s(h)}*W_{8d+2-y},
 \qquad \operatorname{supp} A_h\subseteq\{0,\ldots,L_h\}.
\end{equation}
Zde $A_h$ je nezáporná konečná míra, kterou je potřeba sestrojit.
Ocas má ve všech uzlech stejné hloubky stejnou hmotnost
\begin{equation}\label{T}
 T_d=\E\frac{2}{8d+2-Y_d},\qquad |A_h|=1-T_d.
\end{equation}
Pro $d=0$ je $T_0=1$, $A_\varnothing=0$ a \eqref{normal} dává přesně $W_2$.

\begin{theorem}[Skutečný skalární dyadický odhad]\label{dyadic}
Pro všechna celá $d\ge1$ platí
\begin{align}
 \frac{2}{6d+2}\ \le\ T_d
 &\le\frac{2}{6d+2}
       +\frac{3d}{(6d+2)^2(4d+2)},\label{Tbounds}\\
 \boxed{T_{2d}\le\tfrac12T_d+\tfrac{19}{36}T_d^2.}\label{Tdouble}
\end{align}
Zvláště $T_d=1/(3d)+O(d^{-2})$.
\end{theorem}
\begin{proof}
Položme $Z=Y_d-2d$, $D=6d+2$. Máme
$\E Z=0$, $\E Z^2=3d/2$, $|Z|\le2d$. Použijeme přesnou identitu
\[
 \frac{1}{D-Z}=\frac1D+\frac{Z}{D^2}
                         +\frac{Z^2}{D^2(D-Z)}.
\]
Po zprůměrování a vynásobení dvěma vznikne \eqref{Tbounds}.
Označme nezápornou chybu za $2/(6d+2)$ symbolem $e_d$.
Pak $e_{2d}\le1/(192d^2)$ a
\[
 T_{2d}-\tfrac12T_d
 =\frac{2}{(12d+2)(6d+2)}+e_{2d}-\tfrac12e_d
 \le\frac{19}{576d^2}.
\]
Protože $T_d\ge2/(6d+2)\ge1/(4d)$, dostáváme \eqref{Tdouble}.
\end{proof}

Věta \ref{dyadic} nepostuluje půlení. Odvozuje je pro konkrétní ocasový
účet z \eqref{normal}. Dosud však neříká, že $T_d$ je velikost
\emph{veškerého} kritického defektu: část defektu může vzniknout v hlavě.

\section*{3. Přesné nezáporné doplnění fanu}
Nechť zákon $B_h$ již má tvar \eqref{normal}, $|h|=d$, a položme
\[
 \mu_h=Q*B_h,\qquad M=8(d+1)-s(h),\qquad t=T_{d+1}.
\]
Nad $M$ je zdrojový ocas explicitně znám. Jeho rozdělení v poměrech
$P_x$ dává přesně ocasy dětí předepsané vzorcem \eqref{normal}.
Zbývá matice $\pi_x(u)$ pro $0\le u\le M$, $0\le x\le4$:
\begin{align}
 &\pi_x(u)\ge0,\qquad \pi_x(u)=0\quad(u<x),\label{pos}\\
 &\sum_x\pi_x(u)=\mu_h(u),\qquad
   \sum_{u=0}^M\pi_x(u)=P_x(1-t).\label{rows}
\end{align}
Konečné hlavy potomků jsou pak přesně
\begin{equation}\label{head}
 A_{hx}(j)=\frac{\pi_x(j+x)}{P_x}\quad(0\le j\le M-x).
\end{equation}

Zdroj pro tuto konečnou matici lze navíc napsat zcela explicitně jako
$Q*A_h+D_h$. Položíme-li $a_y=8d+2-y$, potom
\begin{equation}\label{birth}
 D_h(L_h+v)=
 2\sum_{y=0}^{4d}q_d(y)
   \sum_{z=0}^{\min(4,v-1)}
 \frac{Q_z}{(a_y+v-z-1)(a_y+v-z)},\qquad 1\le v\le8,
\end{equation}
a jinde $D_h=0$. Jsou to právě atomy uvolněné rozbalením jednotlivých
zdrojových fází o $8-z$ míst. Tedy celý neznámý systém má konečnou
afinní rekurzi
\begin{equation}\label{affine}
 \boxed{\sum_xP_x\delta_x*A_{hx}=Q*A_h+D_h,\qquad
 |A_{hx}|=1-T_{d+1},\quad A_{hx}\ge0.}
\end{equation}
Ocasový zdroj $D_h$ je znám před volbou fanu a je společný všem jeho
dětem; nevytváříme pro různé budoucí větve různé donory.

\begin{theorem}[Jednokrokový orbitální fan]\label{fan}
Pro daný $B_h$ existují nezáporné hlavy dětí tvaru \eqref{normal}
splňující exactní identitu
$Q*B_h=\sum_xP_x\delta_x*B_{hx}$ právě tehdy, když
\begin{equation}\label{hall}
 \boxed{F_{Q*B_h}(k)\le(1-T_{d+1})F_P(k),\qquad k=0,1,2,3.}
\end{equation}
Při platnosti \eqref{hall} lze hlavy explicitně sestrojit monotónním
transportem konečných měr v \eqref{rows}.
\end{theorem}
\begin{proof}
Zdroj nad $M$ má stejnou hmotnost $t$ jako předepsaný součet ocasů dětí.
Konečný zdroj má hmotnost $1-t$. Podmínka $x\le u$ pro coupling konečného
zdroje a řádkové míry $(1-t)P$ je ekvivalentní kumulativním nerovnostem.
Pro $k=0,1,2,3$ jsou to přesně \eqref{hall}; pro $k\ge4$ je řádková
kumulativní hmotnost již $1-t$, takže nerovnost platí automaticky.
Monotónní kvantilový transport ji realizuje. Vzorec \eqref{head}
a již pevně rozdělený ocas dávají exactnost koeficient po koeficientu.
Vstup lze nejprve vylosovat čerstvě podle $Q$ a symbol vybrat podmíněně
podle této transportní matice; proto jde o causal fan.
\end{proof}

Tato věta pokrývá libovolnou orbitální hloubku a nevyžaduje společný
packet. Sama podmínka \eqref{hall} se ale zatím neukázala jako invariantní.

\section*{4. První přesná nerovnost, kterou ještě neumíme iterovat}
Při výběru současné matice lze bezpečnost všech dětí pro příští fan
zapsat bez jejich nekonečných ocasů. Jejich ocas začíná až za indexem
$M-x\ge4d+4$, takže pro $k\le3$ do něj konvoluce s $Q$ nevidí.
Vedle \eqref{pos}--\eqref{rows} proto potřebujeme současně
\begin{equation}\label{next}
 \boxed{
 \sum_{u=x}^{x+k}F_Q(k+x-u)\pi_x(u)
 \le P_x(1-T_{d+2})F_P(k),
 \quad x=0,\ldots,4,\quad k=0,\ldots,3.}
\end{equation}
Přesně tyto nerovnosti zajišťují \eqref{hall} ve všech pěti dětech.
Jejich jednorázové splnění ještě nezajišťuje, že lze tentýž požadavek
splnit také při výběru všech následujících matic.

\noindent\textbf{Neprokázaný krok této konstrukce:}
existence volby $\pi_h$ ze zdroje $W_2$, která splňuje
\eqref{pos}--\eqref{rows} a \eqref{next} rekurzivně pro všechny historie.
Ocasová část, její fázová orbita i skalární dyadický odhad jsou již
explicitní. Rekurzivní pozitivita hlav není odvozena z půlení ocasové hmoty.
Jde o konkrétní konstrukční ansatz; jeho případné selhání by nevyvracelo
obecnou logaritmickou horní mez.

\section*{5. Přesný účet konečných hlav, pokud orbitální fan existuje}
Uvolněná hmotnost $\Delta_d=|D_h|=T_d-T_{d+1}$ nezávisí na historii.
Z \eqref{peel} plyne
\[
 \Delta_d=\E\frac{2(8-Y)}{a(a+8-Y)},\qquad
 a=8d+2-Y_d,\quad Y\sim Q\text{ nezávisle na }Y_d.
\]
Proto pro $d\ge1$
\[
 \frac{8}{(8d+2)(8d+10)}\le\Delta_d
 \le\frac{16}{(4d+2)(4d+6)}.
\]
Hmotnost $\Theta(d^{-2})$ se tedy uvolňuje ve vzdálenosti
$L_h+1,\ldots,L_h+8=\Theta(d)$ pro \emph{každou} historii.
Její rezervní účet je $\Theta(1/d)$. To je explicitní harmonický
zdroj; nevyvozuje se z požadovaného logaritmu.

Z \eqref{affine}, rovnosti $\E P=\E Q=2$ a předepsaných hmotností hlav
navíc přímo dostáváme
\begin{equation}\label{headincrement}
 \boxed{
 \sum_xP_x\sum_j j A_{hx}(j)-\sum_j j A_h(j)
 =\sum_u(u-2)D_h(u).}
\end{equation}
Pro $d\ge1$ je pravá strana kladná a nejvýše $16/d$.
Pokud se podaří realizovat nezáporné fany do hloubky $n$, součet těchto
přírůstků je tudíž $O(\log n)$. Tato část účtu je již vyřešena; podmínkou
zůstává realizovatelnost nezáporných hlav.

Lze také vyčíslit přesný průměrný účet. Pro celočíselné fáze máme
\[
 \E\min(r+W_a,N)=a\HH_N+r+1-a\HH_a+o(1).
\]
Každý zákon \eqref{normal} má tedy koeficient $2$ před $\HH_N$.
Označme konstantní člen $c_h$. Z exactnosti fanu a rovnosti
$\E P=\E Q=2$ plyne $c_h=\sum_xP_xc_{hx}$. V kořeni je $c_\varnothing=-2$.
Po sečtení přes historie hloubky $d$ dostáváme přesný účet
\begin{equation}\label{meanhead}
 \boxed{
 \sum_{|h|=d}p(h)\sum_j j A_h(j)
 =2\E\HH_{8d+2-Y_d}-(6d+1)T_d-2
 =2\log d+O(1).}
\end{equation}
Při odvození jsme použili $\sum_hp(h)s(h)=2d$ a explicitní ocas
\eqref{normal}; pro každý pevný $d$ jde pouze o konečné součty.
Vzorec \eqref{meanhead} je účetní důsledek existence příslušného konečného
exactního stromu. Není důkazem existence takového stromu.

\medskip
\noindent\textbf{Závěr.}
Orbitální ocasy \eqref{lambda}--\eqref{normal} jsou uzavřeny přesně a
\eqref{Tdouble} je dokázaná all-scale nerovnost. Úplná orbitální
konstrukce vyžaduje ještě rekurzivní pozitivní volbu hlav v
\eqref{pos}--\eqref{next}. Tento dokument tento poslední krok neřeší.
Žádná uniformní těsnost přes historie, omezený posun, omezená fáze ani
kompatibilita nezávisle vybraných konečných svědků se nepředpokládá.
\end{document}
